% GENERATED FILE. Edit canonical lessons, then run npm run generate:books.
% canonical-source-hash: fnv1a64:cac510a072ba807e
% canonical-lessons: PA-C360-imtihan, PA-C360-pepar, PA-C360-ank, PA-C360-natija, PA-C360-pas

\chapter{Exam Time}
\label{ch:pa-a2-exam-time}
\hlchaptermodality{360}

\begin{chapteropening}
\textbf{By the end of this chapter:} \emph{I can talk about my exams and my results.}

Complete the last lesson of chapter 360: I can talk about my exams and my results.
\end{chapteropening}

\section[\texorpdfstring{\pa{ਇਮਤਿਹਾਨ} (imtihān)}{ਇਮਤਿਹਾਨ (imtihān)}]{\pa{ਇਮਤਿਹਾਨ} (imtihān) — an examination; a test}

\label{lesson:PA-C360-imtihan}



\begin{quote}
Before the new one: say the Punjabi for attention, then the Punjabi for to put right.
\end{quote}

\subsection*{You'll want to know: \pa{ਇਮਤਿਹਾਨ}}
\textbf{\pa{ਇਮਤਿਹਾਨ}} — \emph{imtihān} — \textquotedblleft{}an examination; a test\textquotedblright{}.

\begin{quote}
\pa{ਅਗਲੇ} \pa{ਹਫ਼ਤੇ} \pa{ਸਾਡਾ} \pa{ਇਮਤਿਹਾਨ} \pa{ਹੈ।}
\end{quote}

\begin{quote}
\emph{agle hafte sāḍā imtihān hai.}
\end{quote}

\begin{quote}
We have an exam next week.
\end{quote}

A masculine noun.

\begin{cousinweb}
\textbf{\pa{ਇਮਤਿਹਾਨ}} (\emph{imtihān}) is a loan, usually traced through Persian to Arabic \emph{imtiḥān}, \textquotedblleft{}examination\textquotedblright{}.
\end{cousinweb}

\begin{grammarlens}[title={What you've built}]
One more word for this chapter.
\end{grammarlens}

\subsection*{Your turn}
\begin{itemize}
\raggedright
  \item \emph{Say it:} \emph{imtihān}
  \item \emph{Say it:} \emph{imtihān}, once more
  \item \emph{Say it:} \emph{agle hafte sāḍā imtihān hai.} — We have an exam next week.
  \item \emph{Recall:} say the Punjabi for attention, then the Punjabi for to put right, then say \emph{imtihān} again
\end{itemize}

\begin{tcolorbox}[breakable,colback=teal!4,colframe=teal!35!black,title={Before you move on}]
What does \pa{ਇਮਤਿਹਾਨ} mean? (\textquotedblleft{}An examination; a test\textquotedblright{}.) Say it once more.
\end{tcolorbox}
\section[\texorpdfstring{\pa{ਪੇਪਰ} (pepar)}{ਪੇਪਰ (pepar)}]{\pa{ਪੇਪਰ} (pepar) — an exam paper}

\label{lesson:PA-C360-pepar}



\begin{quote}
Before the new one: say the Punjabi for to put right, then the Punjabi for an examination.
\end{quote}

\subsection*{You'll want to know: \pa{ਪੇਪਰ}}
\textbf{\pa{ਪੇਪਰ}} — \emph{pepar} — \textquotedblleft{}an exam paper\textquotedblright{}.

\begin{quote}
\pa{ਅੱਜ} \pa{ਦਾ} \pa{ਪੇਪਰ} \pa{ਔਖਾ} \pa{ਸੀ।}
\end{quote}

\begin{quote}
\emph{ajj dā pepar aukhā sī.}
\end{quote}

\begin{quote}
Today's paper was hard.
\end{quote}

A masculine noun.

\begin{cousinweb}
\textbf{\pa{ਪੇਪਰ}} (\emph{pepar}) is borrowed from English \emph{paper}.
\end{cousinweb}

\begin{grammarlens}[title={What you've built}]
One more word for this chapter.
\end{grammarlens}

\subsection*{Your turn}
\begin{itemize}
\raggedright
  \item \emph{Say it:} \emph{pepar}
  \item \emph{Say it:} \emph{pepar}, once more
  \item \emph{Say it:} \emph{ajj dā pepar aukhā sī.} — Today's paper was hard.
  \item \emph{Recall:} say the Punjabi for to put right, then the Punjabi for an examination, then say \emph{pepar} again
  \item \emph{Write it:} copy \textbf{\pa{ਪੇਪਰ}} once, checking the whole word
\end{itemize}

\begin{tcolorbox}[breakable,colback=teal!4,colframe=teal!35!black,title={Before you move on}]
What does \pa{ਪੇਪਰ} mean? (\textquotedblleft{}An exam paper\textquotedblright{}.) Say it once more.
\end{tcolorbox}
\section[\texorpdfstring{\pa{ਅੰਕ} (aṅk)}{ਅੰਕ (aṅk)}]{\pa{ਅੰਕ} (aṅk) — a mark in a test; a figure}

\label{lesson:PA-C360-ank}



\begin{quote}
Before the new one: say the Punjabi for an examination, then the Punjabi for an exam paper.
\end{quote}

\subsection*{You'll want to know: \pa{ਅੰਕ}}
\textbf{\pa{ਅੰਕ}} — \emph{aṅk} — \textquotedblleft{}a mark in a test; a figure\textquotedblright{}.

\begin{quote}
\pa{ਮੈਨੂੰ} \pa{ਗਣਿਤ} \pa{ਵਿੱਚ} \pa{ਚੰਗੇ} \pa{ਅੰਕ} \pa{ਮਿਲੇ।}
\end{quote}

\begin{quote}
\emph{mainū̃ gaṇit vich chaṅge aṅk mile.}
\end{quote}

\begin{quote}
I got good marks in mathematics.
\end{quote}

A masculine noun.

\begin{cousinweb}
\textbf{\pa{ਅੰਕ}} (\emph{aṅk}) is taken over from Sanskrit \emph{aṅka}, \textquotedblleft{}mark\textquotedblright{}, a tatsama (Monier-Williams).
\end{cousinweb}

\begin{grammarlens}[title={What you've built}]
One more word for this chapter.
\end{grammarlens}

\subsection*{Your turn}
\begin{itemize}
\raggedright
  \item \emph{Say it:} \emph{aṅk}
  \item \emph{Say it:} \emph{aṅk}, once more
  \item \emph{Say it:} \emph{mainū̃ gaṇit vich chaṅge aṅk mile.} — I got good marks in mathematics.
  \item \emph{Recall:} say the Punjabi for an examination, then the Punjabi for an exam paper, then say \emph{aṅk} again
  \item \emph{Write it:} copy \textbf{\pa{ਅੰਕ}} once, checking the whole word
\end{itemize}

\begin{tcolorbox}[breakable,colback=teal!4,colframe=teal!35!black,title={Before you move on}]
What does \pa{ਅੰਕ} mean? (\textquotedblleft{}A mark in a test; a figure\textquotedblright{}.) Say it once more.
\end{tcolorbox}
\section[\texorpdfstring{\pa{ਨਤੀਜਾ} (natījā)}{ਨਤੀਜਾ (natījā)}]{\pa{ਨਤੀਜਾ} (natījā) — a result}

\label{lesson:PA-C360-natija}



\begin{quote}
Before the new one: say the Punjabi for an exam paper, then the Punjabi for a mark in a test.
\end{quote}

\subsection*{You'll want to know: \pa{ਨਤੀਜਾ}}
\textbf{\pa{ਨਤੀਜਾ}} — \emph{natījā} — \textquotedblleft{}a result\textquotedblright{}.

\begin{quote}
\pa{ਇਮਤਿਹਾਨ} \pa{ਦਾ} \pa{ਨਤੀਜਾ} \pa{ਕੱਲ੍ਹ} \pa{ਆਵੇਗਾ।}
\end{quote}

\begin{quote}
\emph{imtihān dā natījā kall āvegā.}
\end{quote}

\begin{quote}
The exam result will come out tomorrow.
\end{quote}

A masculine noun.

\begin{cousinweb}
\textbf{\pa{ਨਤੀਜਾ}} (\emph{natījā}) is a loan, usually traced through Persian to Arabic \emph{natīja}, \textquotedblleft{}result\textquotedblright{}.
\end{cousinweb}

\begin{grammarlens}[title={What you've built}]
One more word for this chapter.
\end{grammarlens}

\subsection*{Your turn}
\begin{itemize}
\raggedright
  \item \emph{Say it:} \emph{natījā}
  \item \emph{Say it:} \emph{natījā}, once more
  \item \emph{Say it:} \emph{imtihān dā natījā kall āvegā.} — The exam result will come out tomorrow.
  \item \emph{Recall:} say the Punjabi for an exam paper, then the Punjabi for a mark in a test, then say \emph{natījā} again
  \item \emph{Write it:} copy \textbf{\pa{ਨਤੀਜਾ}} once, checking the whole word
\end{itemize}

\begin{tcolorbox}[breakable,colback=teal!4,colframe=teal!35!black,title={Before you move on}]
What does \pa{ਨਤੀਜਾ} mean? (\textquotedblleft{}A result\textquotedblright{}.) Say it once more.
\end{tcolorbox}
\section[\texorpdfstring{\pa{ਪਾਸ} (pās)}{ਪਾਸ (pās)}]{\pa{ਪਾਸ} (pās) — passed (an exam)}

\label{lesson:PA-C360-pas}



\begin{quote}
Before the new one: say the Punjabi for a mark in a test, then the Punjabi for a result.
\end{quote}

\subsection*{You'll want to know: \pa{ਪਾਸ}}
\textbf{\pa{ਪਾਸ}} — \emph{pās} — \textquotedblleft{}passed (an exam)\textquotedblright{}.

\begin{quote}
\pa{ਮੇਰਾ} \pa{ਭਰਾ} \pa{ਇਮਤਿਹਾਨ} \pa{ਵਿੱਚ} \pa{ਪਾਸ} \pa{ਹੋਇਆ।}
\end{quote}

\begin{quote}
\emph{merā bharā imtihān vich pās hoiā.}
\end{quote}

\begin{quote}
My brother passed the exam.
\end{quote}

\begin{cousinweb}
\textbf{\pa{ਪਾਸ}} (\emph{pās}) is borrowed from English \emph{pass}.
\end{cousinweb}

\begin{grammarlens}[title={What you've built}]
One more word for this chapter.
\end{grammarlens}

\subsection*{Your turn}
\begin{itemize}
\raggedright
  \item \emph{Say it:} \emph{pās}
  \item \emph{Say it:} \emph{pās}, once more
  \item \emph{Say it:} \emph{merā bharā imtihān vich pās hoiā.} — My brother passed the exam.
  \item \emph{Recall:} say the Punjabi for a mark in a test, then the Punjabi for a result, then say \emph{pās} again
  \item \emph{Write it:} copy \textbf{\pa{ਪਾਸ}} once, checking the whole word
\end{itemize}

\begin{tcolorbox}[breakable,colback=teal!4,colframe=teal!35!black,title={Before you move on}]
What does \pa{ਪਾਸ} mean? (\textquotedblleft{}Passed (an exam)\textquotedblright{}.) Say it once more.
\end{tcolorbox}
